% An input fragment. The enclosing article supplies amsmath/amssymb/hyperref.
\section{A modular criterion for improving the crossing method}
\label{sec:modular-rank-criterion}

The new research direction is to study the ranks of the completed
Heisenberg crossing pieces, after their locality cancellations, rather
than the dimensions of the intermediate Fock spaces in the displayed
factorizations. The following compiler is exact. The asymptotic rank
estimate that would give a better SAT exponent remains open.

The source is Appendix A, Definition 7 and Theorem A, and Appendix B,
equations (53)--(56), of
\href{https://github.com/anthropics/formal-math/blob/1a6eee293b9e00cb432eef6018a21f523623a9af/eth/paper.pdf}
{Anthropic's manuscript at commit \texttt{1a6eee293b9e00cb432eef6018a21f523623a9af}}.
Its two factorizations have combined index length at most
$2\sqrt2\lambda+O(\log\lambda)$, corresponding to
$\rho=1/\sqrt2$ in its recurrence. These are upper bounds on the
actual matrix ranks. The earlier ProveIt Heisenberg entropy obstruction
concerns changing the graded representation while retaining its
operators. It does not prove that these completed-piece ranks attain
their displayed factor dimensions.

\subsection{Exact rank compilation over a prime field}

\paragraph{Lemma.}
Let $N=2^\lambda$, $\lambda\ge1$. Suppose polynomially many integer
arrays $A_i:\{0,\ldots,N-1\}^4\to\mathbb Z$ satisfy
\begin{equation}\label{eq:modular-crossing-identity}
 \sum_i A_i(x,y,u,v)=d_\lambda[x=v][y=u].
\end{equation}
Their entries have polynomial bit length in $\lambda$ and can be
computed in $2^{O(\lambda)}$ bit operations. For a prime
$p\nmid d_\lambda$, let $r_i(p)$ and $c_i(p)$ be the ranks over
$\mathbb F_p$ of the flattenings
\[
 R_i[(x,y),(u,v)]=A_i(x,y,u,v),\qquad
 C_i[(x,u),(y,v)]=A_i(x,y,u,v).
\]
Then $A_i\bmod p$ has exact factorizations through respectively
\[
 b_i^{\rm row}(p)=\left\lceil\log_2\max\{1,r_i(p)\}\right\rceil,
 \qquad
 b_i^{\rm col}(p)=\left\lceil\log_2\max\{1,c_i(p)\}\right\rceil
\]
Boolean bits. The factors can be constructed and verified in
$2^{O(\lambda)}\operatorname{poly}(\log p)$ bit operations and space,
with every scalar represented by an integer in $[0,p-1]$.

\paragraph{Proof.}
Each flattening has order $D=N^2$. Perform Gaussian elimination over
$\mathbb F_p$, obtaining a reduced echelon form of rank $r$. Let $J$
be its pivot columns, let $U$ be the original matrix restricted to
those columns, and let $V$ be the nonzero rows of the echelon form.
The pivot columns of $V$ form the identity. Since row operations
preserve column relations, every original column is the combination
of the columns of $U$ with coefficients in its column of $V$.
Thus $M=UV$, and the columns of $U$ are independent. Pad the inner
index to the next power of two by zero summands. For $r=0$, use one
identically zero summand and no index bits. Apply this construction
to both flattenings.

Elimination uses $O(D^3)=O(2^{6\lambda})$ field operations and
$O(D^2)$ field elements of storage. Field arithmetic has bit cost
polynomial in $\log p$. Generating all entries and multiplying the
factors back both fit $2^{O(\lambda)}\operatorname{poly}(\log p)$.
Reduction modulo $p$ preserves
\eqref{eq:modular-crossing-identity}, and $d_\lambda$ is invertible.
This proves the lemma.\hfill$\square$

\paragraph{The coefficient-height issue.}
Rational elimination does not automatically meet the source's
polynomial bit-length requirement. For $B$-bit entries, a generic
minor determinant of order $D$ has the Hadamard bound
$O(D(B+\log D))=2^{O(\lambda)}$ bits. This is an upper bound,
not a claim that every factorization requires such large coefficients.
It shows that a rational factorization needs an additional height
argument. The modular compiler avoids that obligation: every factor
entry has at most $\lceil\log_2p\rceil$ bits.

The compiler stores exponentially many elements in $\lambda$.
It is not a polynomial-space procedure in that parameter, and no
polynomial-space SAT consequence is claimed here.

\subsection{The stronger uniform rank theorem still needed}

For the source pieces $d_\lambda=h!$, with $h=O(\lambda^2)$. Write
$q_i(\lambda,p)=b_i^{\rm row}(p)+b_i^{\rm col}(p)$. A sufficient
new result would be a fixed
$1/2\le\rho<1/\sqrt2$ and constant $C$ such that
\begin{equation}\label{eq:modular-stronger-rank-goal}
 q_i(\lambda,p)\le4\rho\lambda+C\log_2(2\lambda)
\end{equation}
uniformly for every piece, every parameter used by the recursion,
and sufficiently many primes to reconstruct the answer.
The number of original pieces must remain polynomial in $\lambda$.
The compiler does not replace them by exponentially many pieces.

A uniform bound on the two ranks over $\mathbb Q$ would automatically
establish the same modular upper bound at every prime. Indeed, for an
integer matrix $M$, every minor of order greater than
$\operatorname{rank}_{\mathbb Q}M$ is the zero integer, so
\[
 \operatorname{rank}_{\mathbb F_p}(M\bmod p)
 \le \operatorname{rank}_{\mathbb Q}M.
\]
Thus one may prove a rational-rank theorem but construct its factors
modulo primes, avoiding rational coefficient-height issues. The
converse inference from one favorable prime is invalid.

For example, the concrete rational-rank product target
\[
 \max\{1,\operatorname{rank}_{\mathbb Q}R_i\}
 \max\{1,\operatorname{rank}_{\mathbb Q}C_i\}
 \le 2^{(8/3)\lambda}\lambda^{O(1)}
\]
would give $\rho=2/3$, since the two ceilings add at most two bits.
Its denominator power would be
$\log_2(3/2)=0.5849625\ldots>1/2$. This is a concrete intermediate
target below the source rank-product exponent $2\sqrt2$, and is
also unproved.

\paragraph{The necessary signed-integer bridge.}
Appendix D's truncated parity expansion, Lemma 12 and equations
(64)--(65), is valid modulo powers of two. It cannot simply be
reinterpreted modulo odd primes. Instead, use modular contractions
to reconstruct every requested \emph{integer} clique or network value
exactly, then perform Appendix D's original reductions modulo
$2^\beta$ or $2^{n+1}$ and its exact rational potential comparisons.
This is the integer-value interface in equation (57) and Step 4.

For a Boolean network with $D$ wire bits and vertex entries bounded
by $|W_v|<2^{B_v}$, its integer contraction satisfies
\begin{equation}\label{eq:modular-signed-height}
 |Z|\le2^E,\qquad E=D+\sum_vB_v.
\end{equation}
There are $2^D$ summands, each a product of the indicated entries,
so the triangle inequality proves this bound. A CRT modulus larger
than $2^{E+1}$ therefore recovers $Z$ uniquely as the representative
in $(-P/2,P/2]$. For a $(K,\lambda)$ clique with a uniform table
bound $B$, one may take
\[
 E=\binom K2\lambda+KB.
\]
For the cliques requested by Appendix D Step 4, substitute its
$L\le gP+1$ and $B'=\lceil2gP/K\rceil+g^2B$; these quantities
are polynomial in $n$, so $E=\binom K2L+KB'$ is polynomial as well.
Only each integer value requested by the outer reduction needs
reconstruction. Values internal to its modular contraction need
not be reconstructed individually.

Here is an explicit formulation of the prime requirement. Let $S$
be an input-size parameter at least the encoding length and $n+2$,
and suppose the source cost analysis provides
polynomial caps $\Lambda(S)$ and $H(S)$ on all $\lambda$ and $h$
used in a run, and a polynomial cap $E_{\max}(S)$ on
\eqref{eq:modular-signed-height} for the integer values requested
by the outer reduction. There must be a polynomial-time constructible list
$\mathcal P_S$ of distinct primes, of polynomial cardinality and
polynomial numerical magnitude in $S$, such that
\begin{equation}\label{eq:modular-crt-primes}
 p>H(S)\quad(p\in\mathcal P_S),\qquad
 \prod_{p\in\mathcal P_S}p>2^{E_{\max}(S)+1},
\end{equation}
and \eqref{eq:modular-stronger-rank-goal} holds for all those primes
and all relevant parameters, with the same constants. The first
condition makes every factorial normalizer invertible. A uniform
bound for all primes in a suitable polynomial range suffices;
standard prime-counting estimates supply enough primes there, and
trial division costs polynomial time. Taking $E_{\max}(S)+2$ primes is a
conservative sufficient choice for the product condition.

This is a uniform mathematical requirement, not an instruction to
precompute every possible $\lambda\le\Lambda(S)$ exhaustively.
The latter could be too expensive. Only the pieces requested during
the recursion are compiled, with their costs charged as entry-routine
costs in the source induction.

\paragraph{Conditional consequence.}
The crossing and network identities are algebraic. They remain valid
over $\mathbb F_p$ when all explicitly divided normalizers are units.
Use the compiled factors in the proof of Appendix A's Theorem A.
Their cost is $2^{O(\lambda)}\operatorname{poly}(\log p)$ and their
scalar size is $O(\log p)$, so the same induction has an additional
polynomial dependence on $\log p$ and changed absolute constants.
Its rank recurrence becomes
\[
 e_{\ell+1}=(\rho+O(2^{-\ell}))e_\ell+O(K_\ell^{-1}),
 \qquad \log_2K_\ell=\Theta(2^\ell).
\]
Consequently
\[
 e_\ell=O((\log K_\ell)^{-\alpha}),\qquad
 \alpha=\log_2(1/\rho)>\tfrac12.
\]
Subject to the source's uniform parameter and arithmetic checks with
these modular entry routines and the signed-integer bridge above,
its reductions would give, for fixed $k$,
\[
 \operatorname{poly}(n,m)
 2^{O_k(n\log\log n/(\log n)^\alpha)},
\]
with the corresponding fixed-density bound omitting $\log\log n$.
Each prime run computes a requested integer clique value modulo $p$.
The modulus in \eqref{eq:modular-crt-primes} gives its unique signed
integer reconstruction. The outer algorithm then performs its
original sums, two-adic reductions and rational comparisons.
Polynomially many prime runs per requested value do not change the
exponent order. The smaller modulus bound $>2^n$ suffices for the
direct clause-network route, whose value is already a nonnegative
count, but not for Appendix D's intermediate signed values.

The compiler and this conditional transfer are distinct from the
missing assertion \eqref{eq:modular-stronger-rank-goal}. That assertion
is not proved here. Favorable data at one size or at one exceptional
prime do not imply it, and no existing Lean development is claimed
to have been generalized by this natural-language analysis.

\subsection{Why the straightforward hybrid cancels uniformly}

This statement concerns the displayed upper-bound certificates;
it is not a lower bound on actual running times. Suppose a packet
step has the range budget
\[
 u\log_2r\le B\log_2s
\]
and the mode-bit certificate $S_{\rm cert}=AM/\sqrt u$, for fixed
positive $A,B$. The currently available Anthropic clique coefficient
has $e(r)\ge a/\sqrt{\log_2r}$ on its admissible clique-size sequence,
for fixed $a>0$. Therefore multiplying these estimates gives
\[
 e(r)S_{\rm cert}
 \ge \frac{aAM}{\sqrt{u\log_2r}}
 \ge \frac{aA}{\sqrt B}\frac{M}{\sqrt{\log_2s}}.
\]
This substitution is uniform even for $r=r(s)$. Independently
selecting a different admissible arity at later depths does not
remove the root contribution. Query and padding costs add
nonnegative terms. Thus optimizing this product of existing
certificates cannot give a smaller logarithmic order. A successful
hybrid must change one of its premises, for example by reducing the
actual mode representation or processing stages jointly.

\subsection{Finite validation and essential negative controls}

The standalone checker verifies 945 matrix factorizations, both
flattenings and the crossing identity on nine cases, four rejected
invalid modulus or normalizer inputs, one exceptional-prime profile,
2,049 nonnegative and 4,097 signed CRT reconstructions, and 50 signed
integer network reconstructions followed by two-adic reduction.
It uses the matrix-unit pieces
\[
 A_a(x,y,u,v)=[x=a][v=a][y=u].
\]
Each has row and column rank $N$, giving the attractive combined
cost $2\lambda$. There are, however, $N=2^\lambda$ pieces. This
example deliberately fails the polynomial-piece requirement and
therefore does not establish an improved SAT bound. It guards
against dropping that material hypothesis.

The matrix $\operatorname{diag}(1,3)$ has rank one modulo three
and rank two modulo two, five, and seven. This control shows why
a single favorable prime does not justify the CRT criterion.
The finite checks validate the compiler and reconstruction; they
do not establish asymptotic ranks for the source's Heisenberg pieces.

A separate negative control has one-wire network value
$(-3)\cdot7+5\cdot(-2)=-31$. Its least nonnegative CRT lift modulo
$5\cdot7\cdot11=385$ is $354$. Reducing that lift modulo $32$ gives
$2$, whereas the true signed value gives $1$. Signed reconstruction
first recovers $-31$, after which the two-adic reduction is correct.
